\chapter{Eglot, LSP, and the diagnostic chain}
\label{ch:eglot-lsp}
\minitoc\label{eglot-lsp-and-the-diagnostic-chain}
\index{Eglot}\index{LSP}\index{Language Server Protocol|see{LSP}}\index{flymake}\index{xref}\index{eglot-ensure}

% ── Opening signpost ──────────────────────────────────────────────
The Language Server Protocol gives Emacs deep language intelligence
through external server processes.  This chapter explains how Eglot
maps LSP capabilities to Emacs interfaces, how \texttt{eglot-ensure}
activates servers on \texttt{prog-mode-hook}, and how the
\texttt{flycheck-eglot} bridge unifies diagnostics across LSP and
non-LSP buffers.
This chapter assumes familiarity with
Chapter~\ref{ch:major-minor-modes}.
Chapter~\ref{ch:eglot-server-programs} continues with advanced server
registration; Chapter~\ref{ch:dape} adds the debugger counterpart.

\section{What LSP does}\label{what-lsp-does}

The Language Server Protocol (LSP) defines a JSON-RPC protocol between an editor client and a language server process. The server understands one language deeply (compiles it on the fly, reads \texttt{package.json} / \texttt{pyproject.toml} / \texttt{Cargo.toml}) and exposes language intelligence over a standard wire format. The editor client maps server responses to its own infrastructure.

Eglot is Emacs' built-in (Emacs 29+) LSP client. For each capability in the LSP spec, eglot maps to a corresponding Emacs interface (see Table~\ref{tab:lsp-capabilities}).

\begin{sidewaystable}[p]
\centering
\caption{LSP capability mapping}
\label{tab:lsp-capabilities}
\begin{tabularx}{\linewidth}{@{}XlX@{}}
\toprule
LSP capability & Emacs interface & User command \\
\midrule
\texttt{textDocument/completion} & \texttt{completion-at-point} & \texttt{C-M-i} / corfu auto \\
\texttt{textDocument/hover} & \texttt{eldoc} & echo area / eldoc-box \\
\texttt{textDocument/publishDiagnostics} & \texttt{flymake} (or flycheck via flycheck-eglot) & Fringe indicators \\
\texttt{textDocument/definition} & \texttt{xref} & \texttt{M-.} \\
\texttt{textDocument/references} & \texttt{xref} & \texttt{M-?} \\
\texttt{textDocument/rename} & \texttt{eglot-rename} & \texttt{C-c e r} \\
\texttt{textDocument/codeAction} & \texttt{eglot-code-actions} & \texttt{C-c e a} \\
\texttt{textDocument/documentSymbol} & \texttt{imenu} & \texttt{M-g i} / imenu-list \\
\texttt{textDocument/formatting} & \texttt{eglot-format} & \texttt{C-c e f} / apheleia \\
\bottomrule
\end{tabularx}
\end{sidewaystable}

\section{\texttt{eglot-ensure} and the hook pattern}\label{eglot-ensure-and-the-hook-pattern}

\texttt{eglot-ensure} is the function that starts (or reuses) a server for the current buffer. Adding it to \texttt{prog-mode-hook} means every programming buffer automatically gets an LSP connection:

\begin{minted}{elisp}
(add-hook 'prog-mode-hook #'eglot-ensure)
\end{minted}

This is the idiomatic pattern. Eglot will only start a server if one is registered for the current \texttt{major-mode} in \texttt{eglot-server-programs}. For unregistered modes it signals a user-friendly error.

\texttt{eglot-sync-connect nil} (set in \texttt{:custom}) means the server starts asynchronously; Emacs does not block waiting for the server to initialise. Completions and diagnostics arrive once the server is ready.

\section{The flycheck-eglot bridge}\label{the-flycheck-eglot-bridge}

Eglot's native diagnostic backend is \textbf{flymake}. Flymake is built-in and good, but this config uses \textbf{flycheck} as the unified diagnostic frontend (because flycheck has richer integrations for grammalecte, shellcheck, etc.).

\texttt{flycheck-eglot} wraps eglot's flymake backend and makes it available as a flycheck checker. With \texttt{flycheck-eglot-exclusive t}, the raw flymake indicators are hidden in eglot-managed buffers so only flycheck overlays appear. The diagnostic display is then consistent across all buffers regardless of whether an LSP server is active.

\begin{minted}{text}
eglot  ──wire──>  flymake backend  (eglot internal)
                       │
               flycheck-eglot  (re-exposes as flycheck checker)
                       │
                  flycheck UI  (fringe icons, doom-modeline counts)
\end{minted}

\section{xref: cross-reference infrastructure}\label{xref-cross-reference-infrastructure}

\texttt{xref} is the built-in cross-reference layer that decouples "find definition" from its implementation. \texttt{M-.} (\texttt{xref-find-definitions}) and \texttt{M-?} (\texttt{xref-find-references}) work the same regardless of whether the backend is eglot (LSP), Emacs Lisp introspection, etags, or grep. Eglot registers itself as an xref backend when it manages a buffer.

\texttt{consult-xref} (from the \texttt{consult} package) replaces the default xref UI with a live-filtered consult buffer, giving the same orderless matching and preview available in \texttt{consult-buffer}.

% ── Mental model summary ──────────────────────────────────────────
\section{Mental model}
Eglot's design philosophy is to map each LSP capability onto a standard Emacs interface rather than invent new ones: completions flow through CAPF, hover information through eldoc, diagnostics through flymake, and definitions through xref.  Placing \texttt{eglot-ensure} on \texttt{prog-mode-hook} starts a language server automatically, but only when an entry exists in \texttt{eglot-server-programs} for the buffer's mode.  Setting \texttt{eglot-sync-connect} to \texttt{nil} makes server startup asynchronous, so Emacs never blocks while a server initialises.

Because this configuration uses flycheck rather than flymake for diagnostics, \texttt{flycheck-eglot} bridges the two: it re-exposes flymake diagnostics as flycheck checkers, giving a single consistent diagnostic UI across all buffers.  On the navigation side, \texttt{xref} decouples ``find definition'' from its backend (Eglot, etags, and grep all feed the same \texttt{M-.} interface, and \texttt{consult-xref} layers live filtering on top.  \(\rightarrow\) See Appendix~\ref{app:lsp-protocol} for the JSON-RPC wire protocol details.
